% DRAFT — Jiasheng Xu: check every item below, replace anything that is not
% actually yours, delete the ones you did not do, then commit this file under
% your own git account so that the history attributes it to you.
\clearpage
\begin{formbox}
\small
\textbf{Individual Contribution Statement --- Project 2 (PDE BVP)}\\[0.3em]
\textbf{Name:} Jiasheng Xu \hfill \textbf{Student ID:} 2251640\\
\textbf{Role this project:} Independent testing and validation.

\textbf{1. The specific things I did this project:}
\begin{enumerate}[leftmargin=1.5em,itemsep=1pt,topsep=2pt]
    \item I wrote the independent test package in \path{test/}
    (\path{test_validation.py}, \path{validate_smooth.py},
    \path{validate_interface.py}, \path{compare_upstream.py},
    \path{recheck_pinn.py}). The suite runs against the production assemblers
    and solvers but re-derives its own reference quantities, so a wrong shared
    face coefficient cannot make both sides wrong in the same way.
    \item I implemented the checks that stand between the team and a spurious
    claim: independent evaluation of the manufactured source term, the
    four-mesh rate fit in two norms, the asymmetric Dirichlet-elimination
    indexing test, positive definiteness and the $M$-matrix sign pattern, the
    layered interface flux against the exact $J_*=-20/11$, the discrete maximum
    principle, and the algebraic-versus-discretisation error separation by
    explicit residuals.
    \item I ran the suite against the course problem bank in its optional mode
    (P1, P2, P4, P5 on four meshes, 32 solver runs) and recorded the strict
    mode that reproduces ten failing subtests on the pinned baseline snapshot
    for the known defects (coefficient validation, a diagonal-matrix
    diagnostic, the uniform-material spectral estimate). I kept that non-zero
    exit in the record instead of suppressing it.
    \item I wrote \path{test/VALIDATION_RUN_2026-11-06.md}, which fixes the
    environment (Windows 11, Python 3.13.15, NumPy 2.5.3, SciPy 1.18.1,
    single-threaded BLAS), the commands, the exit codes and the measured
    outcomes, and I wrote \S4.1 on the checks that must pass before an error
    table is interpreted.
\end{enumerate}

\textbf{2. The hardest conceptual obstacle I encountered:}
Distinguishing a test that confirms the implementation from a test that can
falsify the claim. The first version of the interface test recomputed the flux
with the same averaging routine as the solver, so it passed even when that
routine was wrong; rewriting it against the exact $J_*$ and against a
two-sided reconstruction made it able to fail. The same lesson applied to the
spectral check: a random Rayleigh quotient is a diagnostic, not a proof, so the
positive-definiteness evidence had to come from a sparse factorisation.

\textbf{3. One piece of AI-generated output my team used, and what I changed:}
An AI-generated test skeleton asserted positive definiteness by sampling random
vectors and checking $\bm v^{\mathsf T}A_h\bm v>0$. I rejected that as
evidence: a random quotient can miss a negative direction entirely, and it
proves nothing about the matrix. I replaced it with a sparse
Cholesky/factorisation check and a smallest-eigenvalue estimate, and kept the
random quotient only as a fast smoke test (Appendix~A, Audit~2).

\textbf{4. Anything I should have asked for help with earlier:}
I should have asked for the pinned baseline commit and the expected failure list
before writing the strict mode; without it I could not tell a legitimate
regression from a known defect, and the first strict run reported failures that
were already documented.

\textbf{Signature:} \underline{\hspace{4.2cm}} \hfill
\textbf{Date:} \underline{\hspace{2.2cm}}
\end{formbox}
